\documentclass[10pt]{scrreprt}
\usepackage[a4paper, top=30mm, left=25mm, right=25mm, bottom=30mm]{geometry}
\usepackage[utf8]{inputenc}

\usepackage[Bjornstrup]{fncychap}
%\usepackage{ngerman}
\usepackage{graphicx}
\usepackage{epstopdf}
\usepackage{etoolbox}
\usepackage{enumitem}
\usepackage{url}
\usepackage{numprint}
\usepackage{longtable}
\usepackage{caption}
 \usepackage{rotating}
\usepackage{array}
\usepackage{listings}
\usepackage{amssymb}


\makeatletter
\patchcmd{\@makechapterhead}{\vspace*{50\p@}}{\vspace*{-20\p@}}{}{}
\patchcmd{\@makeschapterhead}{\vspace*{50\p@}}{\vspace*{7\p@}}{}{}
\patchcmd{\DOTIS}{\vskip 40\p@}{\vskip -12\p@} 
\makeatother
  
\captionsetup[figure]{labelfont={sf,bf},textfont={sf}}
\deffootnotemark{[\thefootnotemark]}
\deffootnote{1.5em}{1em}{[\thefootnotemark] }
\setlength{\parindent}{0pt}
\renewcommand{\labelitemi}{ \raisebox{0.3ex}{\small$\blacktriangleright$} }
\renewcommand{\labelitemii}{ \raisebox{0.3ex}{\small$\triangleright$} }
\lstset{language=Java}

\newcommand{\sfbf}[1]{\textbf{\sffamily #1}}
\newcommand{\sfit}[1]{\textit{\sffamily #1}}
\newcommand{\W}{\sfbf{W}}
\newcommand{\ziel}[1]{{\fontsize{9.5}{11}\textsf{/#1/}}}
\newcommand{\ziellabel}{Z}
\newcommand{\muss}{\renewcommand{\labelenumi}{\textbf{\ziel{\ziellabel\numprint{\theenumi}0}}}}
\newcommand{\wunsch}{\renewcommand{\labelenumi}{\textbf{\ziel{\ziellabel\numprint{\theenumi}0W}}}}
\newcommand{\JoglEarth}{\raisebox{-1.2mm}{\includegraphics[scale=0.33]{Logo-Text.eps}} }
\newcommand{\textref}[1]{\mbox{\raisebox{0.1ex}{\small$\rightarrow$ }\textit{#1}}}

\newenvironment{details}[1][6pt]{%
  \parskip#1 \parindent6mm \raggedright%
  \def\item{\par\ignorespaces\hangindent=5mm \hangafter1}}{%
  \par\ignorespaces} 
  

\begin{document}

\thispagestyle{empty}
\sffamily
 
\title{Spezifikation}

\begin{figure}
\begin{flushright}
	\includegraphics[scale=0.4]{uniLogo.eps}
\vspace{2.0 cm}
\end{flushright}
\end{figure}

\begin{center}
\vspace{2.0 cm}
{\LARGE SEP – Wintersemester 2013/14}

\vspace{1.0 cm}
\textbf{{\Huge Spezifikation}}

\vspace{0.4 cm}
\begin{figure}[!htb]
\begin{center}
	%\includegraphics[scale=1.0]{projektLogo.eps}
	\includegraphics[scale=1.5]{Logo-Print.eps}
\end{center}
\end{figure}

\vspace{0.2 cm}
\textbf{{\huge OpenStreetMap: Die Welt in 3D}}

\vspace{1.5 cm}
22.11.2013

\vspace{0.5 cm}
Version: 1.0

\vspace{1.5 cm}
{\Large Projektbetreuer: Peter Barth}

\vspace{1.5 cm}
\begin{tabular}{|c|c|c|}
\hline 
\rule[-1ex]{0pt}{4ex} \textbf{Phase} & \textbf{Verantwortlicher} & \textbf{E-Mail Adresse} \\ 
\hline  \hline
\rule[-1ex]{0pt}{4ex} Pflichtenheft & Gabriele Haas & haasgab@fim.uni-passau.de \\ 
\hline  \hline
\rule[-1ex]{0pt}{4ex} Entwurf & Thomas Eder & ederthom@fim.uni-passau.de \\ 
\hline  \hline
\rule[-1ex]{0pt}{4ex} Spezifikation & Christof Blauberger & blauberg@fim.uni-passau.de \\ 
\hline  \hline
\rule[-1ex]{0pt}{4ex} Implementierung & Fabian Knorr & knorrfab@fim.uni-passau.de \\ 
\hline \hline 
\rule[-1ex]{0pt}{4ex} Testing & Constantin Wenger & wengerco@fim.uni-passau.de \\ 
\hline  \hline
\rule[-1ex]{0pt}{4ex} Präsentation & Sebastian Reichl & reichlse@fim.uni-passau.de \\ 
\hline 
\end{tabular}

\end{center}


\pagebreak
\rmfamily
\tableofcontents

\chapter{Einleitung}

\chapter{JavaDocs}

\chapter{JUnit Tests}

\chapter{Algorithmen}

\chapter{Persistente Daten}

\section{Liste der POIs}

\vspace{1.5 cm}
\begin{tabular}{|c|c|c|}
\hline 
\rule[-1ex]{0pt}{4ex} \textbf{ENUM} & \textbf{Erklärung} & \textbf{Overpass Query} \\ 
\hline  \hline
\rule[-1ex]{0pt}{4ex} RESTAURANT & Liefert Cafes, Restaurants und Biergärten &  \\ 
\hline
\rule[-1ex]{0pt}{4ex} NIGHTLIFE & Liefert Bars, Clubs, Diskotheken, Kinos, Casinos und Solarien &  \\ 
\hline
\rule[-1ex]{0pt}{4ex}  BANK & Liefert Banken und Sparkassen &  \\ 
\hline
\rule[-1ex]{0pt}{4ex} FOOD & Liefert Supermärkte, Bäcker, Metzger und Drogerien &  \\ 
\hline
\rule[-1ex]{0pt}{4ex} SHOPS & Liefert Geschäfte aller Art. &  \\ 
\hline
\rule[-1ex]{0pt}{4ex} ACTIVITY & Liefert Freizeitaktivitäten, Picknickplätze und Aussichtspunkte &  \\ 
\hline
\rule[-1ex]{0pt}{4ex} EDUCATION & Liefert Bildungseinrichtungen &  \\ 
\hline
\rule[-1ex]{0pt}{4ex} HEALTH & Liefert Ärzte, Krankenhäuser, Apotheken und Drogerien &  \\ 
\hline
\rule[-1ex]{0pt}{4ex} UTILS &  &  \\ 
\hline
\rule[-1ex]{0pt}{4ex} HOTELS &  &  \\ 
\hline
\end{tabular}


\chapter{Übersetzung der Bedienoberfläche}
\begin{tabular}{|p{7.5cm}|p{7.5cm}|}
\hline 
\textbf{Deutsch} & \textbf{Englisch} \\ 
\hline
\hline 
Ansicht & view \\ 
\hline
Ansichtsmodus & mode \\ 
\hline
Sonnensystem & solarsystem \\
\hline
Satellitenbild & satellite mode \\
\hline
Höhenprofil & heightmap \\
\hline
Für diesen Ort sind keine Informationen verfügbar & No information available to this point. \\
\hline
Punkt markieren & mark point\\
\hline
Längengrad & Longitude \\
\hline
Breitengrad & Latitude \\
\hline
km (Kilometer) & km \\
\hline
Globus & globe \\
\hline
Karte & map \\
\hline
Kinderweltkarte & world map for kids \\
\hline
Ort & place \\
\hline
Suche & search \\
\hline
Ort eingeben & insert place \\
\hline
in der Umgebung & surrounding \\
\hline
Global & global \\
\hline
Suchergebnisse & search results \\
\hline
Markierte Punkte & marked point \\
\hline
Angezeigte Beschriftungen & displaying labels \\
\hline
Markierung entfernen & remove marker \\
\hline
Einstellungen & settings \\
\hline
Sprache & language \\
\hline
Deutsch & german \\
\hline
Englisch & english \\
\hline
Grafikeinstellungen & graphic settings \\
\hline
Antialiasing & Antialiasing \\
\hline
Texturfilter & texture filter \\
\hline
Detaillevel & level of detail \\
\hline
an - aus & on - off \\
\hline
niedrig - mittel - hoch & low - medium - high\\
\hline
Cachegröße & cache size \\
\hline
Arbeitsspeicher & memory \\
\hline
Dateisystem & filesystem \\
\hline
Handbuch & manual \\
\hline
Info & info \\
\hline
Name & name \\
\hline
Name eingeben & insert name \\
\hline
Beschreibung & description \\
\hline
\end{tabular} 
\end{document}